%%%%%%%%%%%%%%%%%%%%%%%%%%%%%%%%%%%%%%%%%%%%%%%%%%%%%%%%%%%%%%%%%%%%%%%%
\section{Gluing and the growth rate}\label{sec:growth}
%%%%%%%%%%%%%%%%%%%%%%%%%%%%%%%%%%%%%%%%%%%%%%%%%%%%%%%%%%%%%%%%%%%%%%%%

\subsection{Terminology}

Let a meander of order $N$ have crossings $1,\ldots,2N$. The \emph{gap} $i$, for $0\le i\le 2N$, is
the stretch of river between crossings $i$ and $i+1$; gaps $0$ and $2N$ are the two ends. An arch
joining crossings $x<y$ \emph{passes over} gap $i$ if $x\le i<y$. An arch is \emph{exterior} if no
other arch on the same side encloses it, and \emph{short} if it joins two neighboring crossings.
The probability $\Pr_N$ and the expectation $\mathbb E_N$ refer to a meander of order $N$ chosen
uniformly at random.

\subsection{Supermultiplicativity and Fekete's lemma}

Two meanders can be glued into one. Let $G$ have order $a$ and $H$ order $b$. Place $H$ to the right
of $G$ on the same river. The last crossing $2a$ of $G$ is the right end of an exterior upper arch
$\{x,2a\}$, and the first crossing of $H$, now numbered $2a+1$, is the left end of an exterior upper
arch $\{2a+1,y\}$. Replace these two arches by $\{x,y\}$ and $\{2a,2a+1\}$. Nothing crosses, because
$\{x,y\}$ passes over everything between $x$ and $y$; and re-pairing one arch of each of two closed
curves joins them into one. Intuitively, both roads are cut open next to the boundary between them
and reconnected across it. The result $G\cdot H$ is a meander of order $a+b$, the \emph{product} of
$G$ and $H$. Given $a$, the pieces are recovered by undoing the re-pairing. Hence
\begin{equation}\label{eq:supermult}
M_{a+b}\;\ge\;M_aM_b\qquad\text{for all }a,b\ge1 .
\end{equation}
Section~\ref{sec:products} identifies exactly which meanders are products.

Inequality~\eqref{eq:supermult} says that $\log M_n$ is superadditive. Fekete's lemma~\cite{Fekete} then says
that $\log M_n/n$ converges to its supremum. The limit is finite because the upper and lower arch
systems are each counted by the Catalan number $\mathrm{Cat}_n=\binom{2n}{n}/(n+1)$, so
$M_n\le\mathrm{Cat}_n^2\le16^n$. Hence $\mu=\lim M_n^{1/n}$
exists, $\mu\le16$, and
\[ M_n\;\le\;\mu^n,\qquad\text{that is}\qquad u_n\le1 \quad\text{and}\quad L(n)\ge0 \qquad\text{for every } n. \]
In terms of $L$, inequality~\eqref{eq:supermult} says that $L$ is \emph{subadditive},
\begin{equation}\label{eq:subadd}
L(a+b)\;\le\;L(a)+L(b),
\end{equation}
and the existence of $\mu$ says $L(n)/n\to0$. These three facts about $L$ (it is nonnegative,
subadditive, and $o(n)$) are used repeatedly.

\subsection{The ratio limit}

\begin{theorem}[ratio limit]\label{thm:ratio}
$M_{n+1}/M_n\to\mu$. In particular the largest ratio $R_\star=\sup_n M_{n+1}/M_n$ is finite.
\end{theorem}

The proof rests on an exact identity. A \emph{wiggle} is a short upper arch followed by a short lower
arch, on three consecutive crossings $i,i+1,i+2$. Any crossing $x$ of a meander of order $n$ can be
replaced by three consecutive crossings that form a wiggle, with the old arches at $x$ reattached to
the outer two. This keeps the road one loop and gives a meander of order $n+1$ with a marked wiggle;
collapsing the marked wiggle undoes it. So marked crossings of order-$n$ meanders correspond exactly to
marked wiggles of order-$(n+1)$ meanders, and with $Z$ the number of wiggles,
\[ M_{n+1}\,\mathbb E_{n+1}[Z] \;=\; 2n\,M_n . \]
A further counting argument with these insertions shows that $Z/n$ concentrates at $2/\mu$ with exponentially small error probabilities, in
both directions. The identity then gives
$M_{n+1}/M_n=2n/\mathbb E_{n+1}[Z]\to\mu$. Since $Z\le n+1$, it also gives $M_{n+1}\ge M_n$: the counts never decrease.

The ratio limit has a rate. Its proof combines a lower bound on $M_n$ with two inequalities that
hold for every $n\ge3$, with a constant $C_0$:
\begin{equation}\label{eq:ratioinput}
\frac{M_{n+1}}{\mu M_n}\;\ge\;1-C_0\sqrt{\frac{L(n)+\log n}{n}},\qquad
\frac{\mu M_n}{M_{n+1}}\;\ge\;1-C_0\sqrt{\frac{L(n+1)+\log n}{n}} .
\end{equation}
The first bounds the ratio from below and the second bounds it from above. Both come from a further
structural decomposition of meanders that we do not describe here~\cite{Repo}. A lower bound on
$M_n$ is an upper bound on the loss $L$, and it makes both right sides close to $1$.

\begin{theorem}[rate in the ratio limit]\label{thm:rate}
There is a constant $C_{\mathrm r}$ such that for every $n\ge2$
\[ \Big|\frac{M_{n+1}}{\mu M_n}-1\Big| \;\le\; C_{\mathrm r}\, n^{-1/4}\sqrt{\log n} . \]
Consequently $|\log(u_{n+k}/u_n)|\le C_{\mathrm r}'\,k\,n^{-1/4}\sqrt{\log n}$ for $0\le k\le n$.
\end{theorem}

The constants are large: with the constants of the proof the bound says something only for $n$
beyond about $10^{14}$. The proof inserts the earlier lower bound $L(n)\le C'\sqrt n\log n$
into~\eqref{eq:ratioinput}. With Theorem~\ref{thm:main} in its place the factor $\sqrt{\log n}$
disappears, by the following short argument, which we give in full.

\begin{remark}
There is a constant $C_{\mathrm r}''$ such that $|M_{n+1}/(\mu M_n)-1|\le C_{\mathrm r}''\,n^{-1/4}$ for every
$n\ge2$. Indeed, by Theorem~\ref{thm:main}, $L(n)\le C\sqrt n$ and $L(n+1)\le2C\sqrt n$, so both right
sides of~\eqref{eq:ratioinput} are at least $1-\epsilon_n$ with $\epsilon_n=O(n^{-1/4})$; when
$\epsilon_n\le1/2$ the second inequality gives $M_{n+1}/(\mu M_n)\le1/(1-\epsilon_n)\le1+2\epsilon_n$, and
the finitely many other $n$ are covered by $M_{n+1}/M_n\le R_\star$ after enlarging the constant.
Summing the logarithms as in Theorem~\ref{thm:rate} gives $|\log(u_{n+k}/u_n)|=O(k\,n^{-1/4})$ for
$0\le k\le n$, and the proof of Proposition~\ref{prop:cutbound} below, run with this rate, improves
the cut bound to $O(s^{-1/4})$. The inequalities~\eqref{eq:ratioinput} cannot give more than
$n^{-1/4}$: sequences of the same form as the decomposition behind them, with loss of order $\sqrt n$, can
have ratios that differ from $\mu$ by a multiple of $n^{-1/4}$ infinitely often~\cite{Repo}.
\end{remark}

Theorem~\ref{thm:rate} says that $u_n$ is \emph{locally smooth}: it changes by a bounded factor over
windows of length about $n^{1/4}/\sqrt{\log n}$, and by the remark over windows of length about
$n^{1/4}$. It does not say that $u_n$ is decreasing: the
statement $M_{n+1}\le\mu M_n$ for all $n$ is open. It holds for every $n\le31$, since the ratios
$M_{n+1}/M_n$ increase in that range and $M_{32}/M_{31}=11.023$ is below the lower bound
$12.2191$ for $\mu$.

\subsection{The halving identity}

The \emph{doubling defect} is
\[ D(m)\;=\;\log\frac{M_{2m}}{M_m^2}\;=\;2L(m)-L(2m) . \]
By~\eqref{eq:supermult}, $D(m)\ge0$. Section~\ref{sec:products} shows that $e^{-D(m)}$ is exactly
the probability that a random meander of order $2m$ is a product of two meanders of order $m$.

\begin{proposition}[halving identity]\label{prop:halving}
For every $n\ge1$,
\[ L(n)\;=\;\sum_{j\ge0}2^{-(j+1)}D(2^jn). \]
\end{proposition}

\begin{proof}
Since $D(m)=2L(m)-L(2m)$, the partial sums telescope:
$\sum_{j=0}^{J-1}2^{-(j+1)}D(2^jn)=L(n)-2^{-J}L(2^Jn)$. The last term is
$n\cdot L(2^Jn)/(2^Jn)$, which tends to $0$ because $L(k)/k\to0$.
\end{proof}

The right side involves only the counts, not $\mu$. Intuitively, the identity starts from a huge
meander, halves it again and again down to order $n$, and adds up what each halving costs. It turns statements about $u_n$ into
statements about the single sequence $D(m)$.

\begin{corollary}\label{cor:halving}
For a constant $\alpha$:
\begin{enumerate}
\item $L(n)=\alpha\log n+O(1)$ for all $n$ if and only if $D(m)=\alpha\log m+O(1)$;
\item $L(n)=\alpha\log n+o(\log n)$ if and only if $D(m)=\alpha\log m+o(\log m)$;
\item $L(n)\le\chi\log n+O(1)$ for some $\chi$ if and only if $D(m)\le\chi'\log m+O(1)$ for some $\chi'$.
\end{enumerate}
\end{corollary}

\begin{proof}
In each part, one direction follows from $D(m)=2L(m)-L(2m)$ and the other from the identity. For
example, in part 1, if $|D(m)-\alpha\log m|\le c$ for all $m$, then
\[ L(n)=\sum_{j\ge0}2^{-(j+1)}\big(\alpha\log n+\alpha j\log 2+O(c)\big)
 =\alpha\log n+\alpha\log2+O(c), \]
using $\sum_j 2^{-(j+1)}=1$ and $\sum_j j\,2^{-(j+1)}=1$. Part 3 is the same computation with inequalities, using $L\ge0$, and part 2 uses that $\epsilon\log(2^jn)\le\epsilon(\log n+j\log2)$.
\end{proof}

So the power law~\eqref{eq:law} is equivalent to $D(m)=\alpha\log m+O(1)$: the probability that a
meander of order $2m$ splits into two meanders of order $m$ should decay like $m^{-\alpha}$. If $u_n\sim An^{-\alpha}$,
then $D(2m)-D(m)$ tends to $\alpha\log 2\approx2.371$; if in addition the correction is smooth, for example
$u_n=An^{-\alpha}(1+O(1/n))$, then $m\,(D(m+1)-D(m))$ tends to $\alpha$.
Table~\ref{tab:D} in Appendix~\ref{app:data} lists exact values. Along the computed range $D$ is
increasing and concave, and both of these quantities increase, reaching $3.13$ at $m=15$ and $2.19$
at $m=8$. Theorem~\ref{thm:main} gives $D(m)\le 2L(m)\le 2C\sqrt m$.
